\documentclass[11pt]{article}
\usepackage[spanish,es-noquoting]{babel}
\usepackage[utf8]{inputenc}
\usepackage[T1]{fontenc}
\usepackage{times}
\usepackage{graphicx}
\usepackage{geometry}
\geometry{
 a4paper,
 total={170mm,257mm},
 left=20mm,
 top=20mm,
 right=25mm,
 bottom=20mm
}

\usepackage{titling}
\usepackage{titlesec}
\titleformat{\section}
  {\normalfont\large\bfseries}{\thesection.}{0.6em}{}
\titlespacing*{\section}{0pt}{18pt}{8pt}

% ---------------------------------------------------------------------
% Bibliografia: mismo mecanismo que plan-paisaje-sonoro/untref-plan.cls con
% la opcion [biblatex] (biblatex-apa, APA 7, backend biber).
%   \textcite{clave}  -> cita narrativa    Autor (2019)
%   \parencite{clave} -> cita parentetica  (Autor, 2019)
%   compilar: pdflatex -> biber -> pdflatex -> pdflatex  (o latexmk)
% ---------------------------------------------------------------------
\usepackage{ifthen}
\usepackage{etoolbox}
\usepackage[style=apa,backend=biber]{biblatex}
\DeclareLanguageMapping{spanish}{spanish-apa}
\setlength{\bibhang}{0.751cm}
\setlength{\bibitemsep}{6pt}
\appto{\biburlsetup}{\Urlmuskip=0mu\relax}
% biblatex-apa usa "&" antes del ultimo autor (forma inglesa); en espanol
% APA usa "y". Mismo parche que untref-plan.cls.
\DeclareDelimFormat[bib,biblist]{finalnamedelim}{%
  \ifthenelse{\value{listcount}>\maxprtauth}
    {}
    {\ifthenelse{\value{liststop}>2}
       {\finalandcomma\addspace\bibstring{and}\space}
       {\addspace\bibstring{and}\space}}}
\DeclareDelimFormat[parencite]{finalnamedelim}%
  {\ifnum\value{liststop}>2 \finalandcomma\fi\addspace\bibstring{and}\space}
\DeclareDelimFormat[fullcite,fullcitebib]{finalnamedelim}%
  {\ifnum\value{liststop}>2 \finalandcomma\fi\addspace\bibstring{and}\space}
\addbibresource{referencias.bib}

\renewcommand\maketitlehooka{%
  \setlength\parindent{0pt}%
  \begin{minipage}{\textwidth}
    \raggedleft
    \emph{Paisaje Sonoro en CABA}
    \null\hfill\emph{Septiembre de 2026, Buenos Aires, Argentina}
  \end{minipage}\vskip 2ex
  \par\vspace{.3cm}
}

\setlength{\parindent}{0.63cm}
\setlength{\parskip}{0pt}

\begin{document}

\title{\textbf{\Large{Paisaje sonoro en los espacios verdes públicos de la Ciudad de Buenos Aires}}}
\author{\textbf{Téc. Abel M. Lamas}\vspace{1ex}}
\date{\normalsize{Ingeniería de Sonido, Universidad Nacional de Tres de Febrero}\\
\textit{abellamasm@gmail.com}\\
\textit{Presentado para su consideración al legislador Emanuel Ferrario}\vspace{-1ex}}

\maketitle

La Ciudad Autónoma de Buenos Aires tiene, según el último Censo, 3.119.304 habitantes \parencite{indec2023}. La superficie de espacio verde público que administra el Gobierno de la Ciudad equivale a 7,03~m² por habitante, pero esa cifra incluye canteros, veredas y cementerios; si se cuentan sólo los parques, las plazas y las plazoletas donde la gente efectivamente pasa tiempo, la relación cae a 3,99~m² por persona \parencite{dgeyc2023}. A eso hay que sumarle que ese número se calcula sobre la población que vive en la Ciudad, sin contar a quienes trabajan o estudian acá y la usan todos los días, de modo que el espacio disponible por persona presente en cualquier momento es, en la práctica, todavía menor.

La Organización Mundial de la Salud recomienda que toda persona tenga un espacio verde de al menos media hectárea a no más de 300 metros de su vivienda \parencite{who2016}. Hoy, más de 350.000 porteños viven lejos de una plaza \parencite{gadea2021}. El propio Código Urbanístico de la Ciudad reconoce el problema y define los «Espacios Verdes de Proximidad» como aquellos a los que se puede llegar caminando en cinco minutos, unos 400 metros, con el objetivo explícito de reducir el estrés urbano y amortiguar la contaminación en todas sus formas \parencite[art.~10.12.8]{codigourbanistico}.

En ese último punto, se menciona a la contaminación en todas sus formas, lo cual incluye al ruido, y es ahí donde aparece un vacío que este trabajo busca llenar. La disponibilidad de espacios verdes no es sólo una cuestión de metros cuadrados, hay evidencia internacional sólida de que el acceso a un espacio verde cercano reduce la molestia por ruido y los síntomas de estrés en los barrios más expuestos al tránsito \parencite{gidlof2007}, y de que los ambientes sonoros evaluados positivamente se asocian con mejor salud autopercibida y con una recuperación más rápida del estrés \parencite{aletta2018}. Es decir, un espacio verde protege más cuando su sonido también es el adecuado, no sólo cuando tiene suficientes árboles y superficie.

La propia Constitución porteña ordena, promover tanto la preservación y restauración de la calidad sonora como la preservación e incremento de los espacios verdes \parencite[art.~27]{constitucioncaba}. Sin embargo, hoy no existe una herramienta que permita verificar si esa doble obligación se cumple en un espacio verde concreto. La normativa de contaminación acústica que sí se tiene, la Ley 1540,  reserva la máxima protección sonora únicamente a un lugar de toda la Ciudad, la Reserva Ecológica Costanera Sur \parencite{ley1540,decreto740}, y no interviene en ningún caso una medición real del ambiente sonoro. La Ley 1540 define apenas cuatro niveles de evaluación sonora, y la mitad de ellos regula el ruido que entra a un ambiente interior, ya sea desde el exterior o desde otro recinto; incluso en el nivel restante, que sí fija un límite para el ambiente exterior, ese límite es un valor de emisión en decibeles por zona, no una evaluación de cómo se percibe el sonido en el lugar. Dentro de esto se tienen las Áreas de Sensibilidad Acústica al Ambiente Exterior (ASAE), las cuales se asignan sólo según el uso del suelo en base al Código Urbanístico, sin hacer mediciones acústicas. Una plaza de barrio, por más que cumpla todos los requisitos de superficie y accesibilidad, puede tener un ambiente sonoro degradado y no existe hoy ningún mecanismo para detectarlo, mucho menos para exigir que se corrija.

El riesgo de fondo es que, aun si en el futuro se creara una herramienta para evaluar el ambiente sonoro de los espacios verdes, reproduzca el mismo abordaje clásico del control de ruido, que se limita a medir la energía sonora y no contempla cómo las personas efectivamente perciben esos sonidos. Frente a ese riesgo, esta investigación propone adoptar el marco teórico del paisaje sonoro, estandarizado en la ISO 12913-1, que define al paisaje sonoro como el ambiente acústico tal como una persona lo percibe y lo experimenta en su contexto, de modo que no alcanza con medir cuántos decibeles hay en un lugar, sino que además hace falta entender cómo esos sonidos son interpretados por quien los escucha \parencite{iso129131}. Para eso, la ISO 12913-2 propone herramientas de evaluación psicoacústica que van más allá del nivel sonoro continuo equivalente en dBA, que es el que se utiliza en la mayoría de los casos para evaluar el ruido. A esas herramientas se suman investigaciones más recientes y avanzadas, como los modelos matemáticos propuestos por el investigador japonés Yoichi Ando, las funciones de autocorrelación (ACF) y de correlación cruzada interaural (IACF), que describen cómo el sistema auditivo humano procesa un sonido y cuentan con respaldo empírico en mediciones de la actividad del nervio auditivo \parencite{ando2001,soeta2015}.

\textcite{hermida2019} adaptó estos modelos de Ando al estudio del paisaje sonoro urbano, calculando esos factores sobre grabaciones tomadas en Bogotá y en Brasilia y comparándolos con los atributos perceptuales que un grupo de oyentes entrenados evaluó en laboratorio a partir de esas mismas grabaciones. Esa línea de trabajo se complementa con otra investigación del mismo autor, realizada en tres parques de Lisboa, que comparó los resultados obtenidos con encuestas hechas en el lugar contra los obtenidos en pruebas de laboratorio \parencite{hermida2017}. Estos estudios son un aporte valioso, pero, como toda investigación científica, sus hallazgos necesitan verificarse con datos relevados directamente en el lugar, con público general y en el contexto de otra ciudad, para sumar evidencia a hipótesis que todavía se están poniendo a prueba. A esto se suma que buena parte de estas herramientas de evaluación por parte de los estandares ISO 12913 fueron pensadas originalmente en inglés, y que la única adaptación validada al castellano que existe hasta ahora se hizo con hablantes de Bogotá \parencite{useche2023}. Aun cuando ya hay antecedentes en la región, todavía falta adaptar ciertos términos al español rioplatense, de modo que los próximos trabajos que se hagan en la región de Buenos Aires y tambien en proximidades puedan apoyarse en ese mismo lineamiento.

Esta investigación propone desarrollar y poner a prueba un método para caracterizar el ambiente sonoro de los espacios verdes de la Ciudad, tomando como caso de estudio un parque público. El objetivo es obtener esas caracterizaciones para aportar evidencia científica que respalde o ponga a prueba las hipótesis planteadas sobre el paisaje sonoro y sobre los modelos matemáticos de Yoichi Ando, de modo que a futuro puedan usarse como sustento en la ejecución, planificación, restauración y preservación de espacios verdes, y resulten útiles a la hora de realizar un estudio de impacto ambiental y acústico. La intención de fondo es que el problema del ruido empiece a pensarse desde herramientas más modernas, y no sólo desde las tradicionales que, si bien siguen siendo útiles hoy, la evidencia ya mostró que resultan insuficientes.

El resultado de este trabajo no reemplaza ninguna política existente, pero puede aportar una base técnica concreta para dos cosas que hoy faltan: un criterio objetivo para evaluar si un espacio verde cumple efectivamente su función de protección frente al ruido, y un antecedente académico argentino en una disciplina, el paisaje sonoro urbano, donde en la región todavía se tienen muy pocos estudios propios. Se comparte este documento con el objetivo de que sirva como iniciativa para pensar, desde el Gobierno de la Ciudad y sus organismos de control, si la calidad sonora de los espacios verdes merece incorporarse como un criterio explícito en la planificación y el mantenimiento del espacio público porteño.

Este documento se entrega al legislador como material de consulta posterior al encuentro sostenido, para que pueda profundizar en tres ejes conversados durante la reunión. El primero es la propuesta de evaluación del paisaje sonoro y la posibilidad de que resulte de utilidad para el trabajo legislativo en la materia. El segundo son las problemáticas vinculadas al ruido en los espacios verdes públicos que identifican los vecinos de la Ciudad, información valiosa tanto para esta investigación como para la agenda legislativa. El tercero es la identificación, junto con el legislador y su equipo, de los parques o plazas que la comunidad considera de mayor relevancia para su vida cotidiana, de modo de orientar hacia alguno de ellos el caso de estudio de esta investigación.

En ese mismo sentido, el acompañamiento del legislador puede resultar especialmente valioso para la convocatoria de participantes del relevamiento in situ, que requiere contar con vecinos dispuestos a recorrer el parque elegido y completar una breve encuesta en el lugar. Se trata de una tarea de reclutamiento laboriosa para un investigador individual, y el vínculo del legislador con organizaciones vecinales, juntas comunales u otros canales de comunicación con la comunidad podría facilitar sensiblemente esa convocatoria.

\section*{Referencias}
\printbibliography[heading=none]

\end{document}
